\section{Hybrid Split: TEE + Obfuscation}
\label{sec:hybrid}

% SCOPE ROLE: highest deployment-readiness family (a §8 finding, NOT assumed). Needs only a
% CPU-TEE (SGX/TDX, ubiquitous) + COMMODITY untrusted GPUs --- no confidential GPU ---
% making it deployable on commodity GPUs, between pure-TEE (\cref{sec:tee}, fast but
% needs CC-GPU) and crypto (\cref{sec:crypto}, no special HW but 100-1000x). Cover only
% the modern practical schemes (GELO, TwinShield, ObfuscaTune, SecureInfer, SCX, Amulet);
% the CNN-era roots (Slalom, DarKnight, ShadowNet, SOTER) are lineage-only prose, and
% TEESlice is cited as the analysis showing why naive (un-obfuscated) split fails.
\dnote{this is the deep family; lead with the commodity-GPU deployability argument}

% --- scope-framing (short role-framing; bodies stay \needswork) ---
Among the surveyed families, the hybrid split offers the most balanced
deployment-readiness profile (\cref{sec:intro})---a ranking we substantiate in
\cref{sec:synthesis}, not assume here. It keeps a small trusted core in a CPU-TEE
(SGX/TDX, ubiquitous) and offloads the heavy linear algebra to an \emph{untrusted
commodity GPU} over obfuscated activations---no confidential GPU required---combining
near-baseline performance with exact fidelity at Tier-2 trust. It thus sits between
pure-TEE (\cref{sec:tee}: near-plaintext, but needs a confidential GPU at scale) and
cryptography (\cref{sec:crypto}: no special hardware, but $100$--$1000\times$). We
cover the modern practical schemes below; the CNN-era roots (Slalom, DarKnight,
ShadowNet, SOTER) are lineage only, and TEESlice is cited as the analysis showing why
a \emph{naive}, un-obfuscated split leaks under public weights.

\subsection{Approach}
% Keep a small trusted core in a TEE; offload the heavy linear algebra to an
% untrusted accelerator over OBFUSCATED activations. Combines TEE's near-plaintext
% speed on the bulk with obfuscation's cheapness, bounded by what the accelerator
% can infer from mixed activations. Map onto §2: TCB = GPU-TEE (small); HbC +
% white-box; observation surface = accelerator activations.
\needswork{describe the split pattern (boundary layers in TEE, bulk GEMM on
untrusted accelerator over obfuscated activations)}

\subsection{Representative Schemes}
The hybrid-split \emph{inference} schemes---GELO (per-batch invertible mixing $U{=}AH$, QKV/O
offloaded), TwinShield (CPU-TEE root + secret sharing for linear \emph{and} non-linear, with
U-Verify integrity), ObfuscaTune ($\sim$5\% of parameters kept in the TEE), SCX (per-session
one-time-key KV-cache, $(\varepsilon,0)$-DP), and Portcullis (TEE-attested PII-anonymization
gateway)---are reported in the shared inference table (\cref{tab:inf-perf}). On-device
model-protection splits (SecureInfer~\cite{secureinfer}, Amulet~\cite{amulet}) defend the
\emph{model} from the device owner and are out of scope (\cref{sec:intro}).
The graph-RAG anonymization scheme PrivGemo (dual-tower KG with HMAC session identifiers)
applies the same anonymization lineage to the \emph{retrieval} layer; it is covered with the
other graph-retrieval schemes in \cref{sec:retrieval}.
\dnote{This family is the substrate the COMPANION paper builds on (GELO
optimizations + a unified static+dynamic design). Here we survey it only.}

\subsection{Performance (Reported)}
\needswork{the offload fraction vs overhead trade; why this family is the
practical frontier for confidential inference}

\subsection{Limitations and Known Attacks}
% GELO names its own leakage: Gram-matrix invariants under orthogonal mixing
% (U^T U = H^T H), mitigated by non-orthogonal A or shield vectors. Security is
% single-batch only.
\needswork{Gram-leak + mitigations \citep{belikov_gelo}; co-located TEE+GPU
requirement; per-batch-only guarantee}
